% !TeX root = ../../Book3.tex
% 纲要标题：### 21.5 消解的极小化与 Betti 数据的解释
% 纲要位置：计划纲要.md，第 2175 行。

\section{消解的极小化与 Betti 数据的解释}
\label{b3:sec:2105}


% 纲要要点（供目录核对）：
%
% * 在标准分次多项式环上识别消解微分中的单位项，构造并消去可缩自由直和项，得到极小消解
% * 以剩余域张量后的零微分识别分次 Betti 数，并与 \(\operatorname{Tor}\) 的分次维数核对
% * 同时记录微分矩阵、齐次次数和矩阵乘积为零的证书；区分正合性证书与仅检查复形条件
% * 比较 Gröbner 退化前后的 Hilbert 数据和消解信息，说明相同 Hilbert 级数不保证相同 Betti 表
%
% ---
%

计算得到的自由消解可能很长，甚至在一个已经可以停止的位置仍留下互相抵消的自由项。极小化的操作十分具体：找到微分矩阵中的非零常数，把它对应的一对基向量单独分出，再删去一个正合的两项复形。关键是同时修改相邻微分，而不能只删矩阵的一行一列。

\subsection{单位微分项、可缩直和项与极小化}
\label{b3:subsec:2105:01}

\begin{proposition}\label{b3:prop:graded-cancellation}
设 \(k\) 为域，\(S=k[x_1,\ldots,x_n]\) 为标准分次多项式环，\(M\) 为有限生成分次 \(S\) 模，\(F_\bullet\to M\) 为各项有限生成且只有有限个非零项的分次自由消解。若某个 \(i\ge1\) 的微分 \(\mathrm{d}_i:F_i\rightarrow F_{i-1}\) 的矩阵含非零常数项，则可作保持次数的可逆换基，把此消解分为另一消解与一个两项复形
\[
0\longrightarrow S(-a)\xrightarrow{1}S(-a)\longrightarrow0
\]
的直和。重复有限次后，剩下的消解极小。
\end{proposition}
\begin{proof}
把所选常数移至左上角并乘以其逆，令相应源基向量为 \(e\)、目标基向量为 \(f\)，其次数均为 \(a\)。用 \(\mathrm{d}_i(e)\) 替换 \(f\)，由于 \(f\) 系数为一，这仍是齐次基；再从其他源基向量减去适当的 \(e\) 倍数，消去它们在该目标向量上的分量。各乘子次数等于所需基次数之差，所以这些初等变换保持次数。于是 \(\mathrm{d}_i\) 分块为 \(1\oplus \mathrm{d}_i'\)。

由 \(\mathrm{d}_{i-1}\mathrm{d}_i=0\)，新目标基向量在 \(\mathrm{d}_{i-1}\) 下为零；由 \(\mathrm{d}_i\mathrm{d}_{i+1}=0\)，\(\mathrm{d}_{i+1}\) 的像在源向量 \(e\) 处的分量为零。这使两项复形真正分裂为直和项。它的恒等微分使其正合，删除它后其余复形仍为消解。每次删除使自由模总秩下降二；总秩有限，故过程停止。停止时全部微分项属于 \(\mathfrak m\)，即为极小消解。
\end{proof}

例如展示 \((x,y,x+y):S(-1)^3\rightarrow S\) 中，第三列是前两列之和。关系 \((-1,-1,1)\) 含单位系数，它与一个冗余生成元组成可缩项。消去后回到 \((x,y)\) 的 Koszul 消解。这说明常数关系为何是计算中首先应寻找的冗余，而不是几何上的新约束。

\subsection{剩余域张量、Tor 与分次 Betti 数}
\label{b3:subsec:2105:02}

对任意分次自由消解 \(F_\bullet\)，
\[
\operatorname{Tor}_i^S(M,k)=H_i(F_\bullet\otimes_Sk).
\]
若微分含单位，张量后相应微分仍非零，故同调维数一般小于该层自由秩；消去可缩项以后，张量后的微分才全部为零，得到式\eqref{b3:eq:graded-betti-tor}。因此从矩阵读出 Betti 数之前，须先验证极小性。

\ProseParagraphBreak
Koszul 复形还提供不先求自由消解的检验方法。变量列是 \(S\) 正则序列，故其 Koszul 复形消解 \(k\)。由\cref{b3:cor:koszul-tor}，以 \(M\) 为系数的 Koszul 同调计算 \(\operatorname{Tor}_i^S(M,k)\)。这里使用\cref{b3:lem:tor-basic-properties}中的平衡性：在双复形 \(F_\bullet\otimes_SK_\bullet\) 中，无论先用 \(F_\bullet\) 消解 \(M\)，还是先用 \(K_\bullet\) 消解 \(k\)，所得总同调相同。这使直接计算 Koszul 同调成为核对 Betti 数的另一种方法。

\ProseParagraphBreak
例如 \(A=k[x,y]/(x,y)^2\) 的最高 Koszul 同调是被 \(x,y\) 同时零化的元素乘 \(e_x\wedge e_y\)。这些元素恰为 \(kx\oplus ky\)，所以 \(\operatorname{Tor}_2^S(A,k)\) 维数为二，内次数均为三，与式\eqref{b3:eq:square-maximal-resolution}的左端完全一致。

\subsection{微分矩阵、齐次次数与正合性证书}
\label{b3:subsec:2105:03}

一个消解的可复查记录应包含基域、变量序、各自由基向量的次数、增广映射以及各微分矩阵。若 \(F_i\) 中第 \(j\) 个基向量次数为 \(a_{ij}\)，则矩阵 \(\mathrm{d}_i\) 的第 \((h,j)\) 项必须齐次且次数为 \(a_{ij}-a_{i-1,h}\)。这项检查能发现很多矩阵虽然乘积为零、却不保持所称分次的错误。

\ProseParagraphBreak
正合性需要另外的证据。可在每一步给出上一微分核的 Gröbner 基、Schreyer 生成证书，以及这些核生成元与当前微分列之间的双向表示。双向表示证明当前像等于全部核，而 \(\mathrm{d}_i\mathrm{d}_{i+1}=0\) 只证明像包含于核。消解最后一项也必须检查单射，不能因为“后面没有再写矩阵”就认定它没有关系。

\ProseParagraphBreak
最简单的反例是把零映射 \(S(-1)\xrightarrow{0}S\) 接到满射 \(S\rightarrow S/(x)\)。两个连续映射的复合为零，但中间核是 \((x)\)，前一像为零，这根本不是消解。相反，若一个有限复形已知正次数同调有限长度，仅核对 Hilbert 级数的交错和也未必能证明每个同调都为零，因为不同层的非零同调可能在交错和中互相抵消。

\ProseParagraphBreak
本节的算法因而可以按以下顺序执行：先求完整核生成元和表示证书，再逐层构造消解，随后通过单位项消去得到极小复形，最后由剩余域张量读 Betti 数。Hilbert 合冲定理保证极小结果的长度上界；它不保证任意未经极小化的求核过程自动在同一层数停止。

\subsection{Gröbner 退化下的 Hilbert 数据与 Betti 表}
\label{b3:subsec:2105:04}

对齐次理想 \(I\subseteq S=k[x_1,\ldots,x_n]\)，\cref{b3:prop:hilbert-initial-count}已经证明初始理想 \(\operatorname{in}(I)\) 与 \(I\) 有相同的分次 Hilbert 函数。计算 Hilbert 函数时，这把多项式关系变成了单项式计数；然而它没有保持全部合冲关系。

\ProseParagraphBreak
在 \(S=k[x,y]\) 中取字典序 \(x\succ y\)，令
\[
I=(x^2,xy+y^2).
\]
\(x^2\) 是非零因子，\(xy+y^2=y(x+y)\) 与 \(x\) 互素，所以在 \(S/(x^2)\) 中也是非零因子：若 \(x^2\mid a\,y(x+y)\)，唯一分解性使 \(x^2\mid a\)。故两个二次式构成正则序列，商的极小消解有 Betti 多项式 \(1,2t^2,t^4\)。

\ProseParagraphBreak
而 Gröbner 基为 \(G=(x^2,xy+y^2,y^3)\)。第三项由
\[
y x^2-x(xy+y^2)+y(xy+y^2)=y^3
\]
得到；第一、第二个元素的 \(S\)-多项式经约化留下 \(y^3\)，第二、第三对留下 \(y^4\)，再由 \(y^3\) 化为零；第一、第三对首项互素。因而 \(\operatorname{in}(I)=(x^2,xy,y^3)\)，正是上一节的 \(J\)。它的极小消解在第一、第二层都多出一个内次数为三的生成元。

\ProseParagraphBreak
所以这个例子有相同的 Hilbert 函数，却有不同的 Betti 数。\subsecref{b3:subsec:1102:04}曾在 \(k[x,y,z]\) 中计算同样两个生成元的商；增加一个自由变量 \(z\) 使 Hilbert 级数乘 \((1-t)^{-1}\)。现在由正则序列的消解得到
\[
H_{k[x,y,z]/(x^2,xy+y^2)}(t)
=\frac{1-2t^2+t^4}{(1-t)^3}
=\frac{1+2t+t^2}{1-t},
\]
与当时的标准单项式计数相符。这次计算还给出了全部生成关系及其次数，是单独一个级数没有记录的信息。由初始理想得到的结果可以帮助猜测原理想的消解，但回到原理想后仍须检验关系及极小性。更一般的 Betti 数比较涉及平坦族和专门化，本章不把这些比较作为已经证明的结论。
